\preludeandfugue{The Varkolak}
\contributor{Jay Rose Ana}

% Key: C minor, with a shift into G minor

\prelude

\subsection*{A Matter of Record.}

Perth, April 1924, a pale-faced madman, never seen since nor before,
ranting an incomprehensible, yet poetic language, over and over again,
was tried by jury to be committed to Claremont Asylum for the
Insane. Uncovered, disguised as a woman and concealing a cleaver and
iron bar upon his personage. A witness close to the case states he was
far into plans to stun by bludgeon, Perth’s Director of Education, and
gouge out his eyes and worse, heinous things, to which the courtroom
gazed on in horrified surprise.

Unable to hold a simple sentence, and yet, replete with notes dated
back through history, the case was examined, and quickly wrapped up in
a sacrilegious mystery.

He hobbled away, uncomfortably on a withered calf and an inward
twisted ankle, no statement on his behalf. Upon a technicality he was
freed and never seen nor heard from again publicly.

When asked for a statement the Director of Education refused and fell
silent for the rest of his days.

The Corso Polaris camera, serial 130, was reported to be devoid of the
image of the accused due to a suspected manufacturing imperfection and
being improperly used.

The camera was seized by the court for further investigation. No
investigation was ever concluded.

\fugue

\subsection*{A Matter of Time.}

One hundred years earlier, April 1824 a dreamer, a darer, a young
poet, excommunicated from the church, contracted an unusual fever
during the Ottoman attempt to capture the strategic town Missolonghi
amidst the Greek War of Independence. He succumbed to fever,
malevolent in its intent. Allegedly, the fever was caused by
sporozoites of malaria parasites that burrowed deep into his liver,
infecting red blood cells and feasting on the liver itself whilst in
the distance gunfire rang forebodingly like distant choking.

Though he had tried, for many years, he was unable to fight this
infection to the end and this dreamer succumbed and dreamed his last
dream in a medical tent, and even sunlight could not darken his
complexion, say the physicians Bruno and Millingen as they worked
feverishly to cut the damaged liver away, though both seemed very
young and inexperienced for the day.

In 1817, Venice, shortly after the prohibited Venice Carnival, there
is a fable, a tale, La Repubblica Serenissima, of an opulent feast, an
impossible, outlandish banquet, floating upon pontoons in the lagoon
of Rivo Altus, between Santa Croce and Punto San Giuliano.

In attendance, purportedly, if witnesses were to be believed were Pope
Gregory, Marco Polo, Henry Tudor, King Louis XIV, Giacomo Casanova,
and our young dreamer poet travelling Europe with a limp. Escorting
this young dreamer were three physicians, Polidori, Bruno, and
Millingen. Of the three, Polidori had literary ambitions, penning a
fantastical story The Vampyre in 1816, whilst in the company of
Shelley and our friend, and shortly after was dismissed, returning to
Norwich, England, where he allegedly committed suicide in 1821.

Or so goes the story recounted by a certain someone who tells a
slightly different story in person.

In 1980, The Carnival of Venice resumed and has run every year since.

The opening day of the Carnival of Venice on 27 January 2024, where
one might have their bauta, tricorn, and black zendale, ready for who
knows what fantabulous tale they may have to endure was veiled in
thick fog. In one of the news photographs, from that very Carnival,
people dressed as La Bauta could be seen, set against heavy Venetian
fog, and it is reported one of those individuals expressed a very
heavy limp.

This year, 2026, saw Olympus, The Origins of the Game, with the
opening traditional event taking place on water, a procession along
the Grand Canal, culminating beneath the Rialto Bridge. When asked of
the event, and of what happened beneath that bridge, the crowd fell
unusually silent and their cameras unusually devoid of images.

\endinput

Sadie, thank you for asking me about this. I thought about this piece
which I wrote for a poetry evening a few years ago, I read it there
once, but it is unpublished. It is based on a mix of true events and
fiction and I have reworked and updated it for your format,

If you can let me know if it hits the spot of what you are looking for
and anything you feel might need to change.

best wishes

Jay Rose Ana

Title: The Varkolak

Author: Jay Rose Ana

Prelude: A Matter of Record

Fugue: A Matter of Time

Key: C minor, with a shift into G minor

Publication status: Unpublished. Previously performed once at a poetry
evening several years ago.

Preferred pronouns: she/her

Email: jayroseana@gmail.com

Postal address: This is my family home address although I typically
live in my flat in Worcester (this address is more reliable)

7 Priory Way
St Georges
Telford
Shropshire
TF2 9YQ

Biography:

Jay Rose Ana is a poet, spoken word performer and the current
Worcestershire Poet Laureate. She writes about people, places, the
things we carry, and the wonderfully ordinary moments that somehow
become poems. She believes poetry belongs everywhere, and can usually
be found with a notebook, a cappuccino and a lovely piece of cake. She
is currently spending her Laureate year listening, writing, making
things happen and, wherever possible, being fabulous.

Agreement to the rules:

I confirm that I have read and agree to the rules for the project.


The Varkolak

Jay Rose Ana

Key: C minor, with a shift into G minor

Prelude.

A Matter of Record



Fugue.

A Matter of Time

 

 
